\documentclass[reprint,aps,prl,twocolumn,superscriptaddress]{revtex4-1}
\usepackage{braket}
\usepackage{amssymb}
\usepackage{amsmath}
\usepackage{epsfig}
\usepackage{color}
\usepackage{graphics, graphicx}
\usepackage{bbold}
\usepackage{psfrag}
\usepackage{mathcomp}
\usepackage{subfigure}
\usepackage{verbatim}
\usepackage[colorlinks, citecolor=blue]{hyperref}
\usepackage[normalem]{ulem}
\usepackage{bm}
\def\cp#1{\mathbf{#1}}
\begin{document}


\title{Self-bound Vortex Lattice in a Rapidly Rotating Quantum Droplet}
\author{Qi Gu}
\affiliation{Institute for Advanced Study, Tsinghua University, Beijing 100084, China}
\author{Xiaoling Cui}
\email{xlcui@iphy.ac.cn}
\affiliation{Beijing National Laboratory for Condensed Matter Physics, Institute of Physics, Chinese Academy of Sciences, Beijing 100190, China}

\date{\today} 


\begin{abstract}
A rapidly rotating Bose gas in the quantum Hall limit is usually associated with a melted vortex lattice. 
In this work, we report a self-bound and visible triangular vortex lattice without melting for a two-dimensional Bose-Bose droplet rotating in the quantum Hall limit, i.e., with rotation frequency $\Omega$ approaching the trapping frequency $\omega$.
Increasing $\Omega$ with respect to interaction strength $U$, we find a smooth crossover of vortex lattice droplet from a {\it needling} regime, as featured by small vortex cores and an equilibrium flat-top surface, to the lowest-Landau-level regime with Gaussian-extended cores spreading over the whole surface.   
The surface density of such rotating droplet is higher than that of a static one, and their ratio is found to be a universal function of $\Omega/U$. We have demonstrated these results by both numerical and variational methods. The results pave the way for future experimental exploration of rapidly rotating ultracold droplets into the quantum Hall limit.  
\end{abstract}
\maketitle


Ultracold Bose gases in rotating harmonic traps have attracted much attention in recent years\cite{review1,review2}. An important motivation is the emergence of quantum Hall physics when the rotating frequency($\Omega$) approaches the trapping frequency($\omega$)\cite{JasonHo, Cooper}. This situation is analogous to charged particles in a magnetic field, and a triangular vortex lattice emerges therein similar to  those  in type-II superconductors\cite{Abrikosov1957} and rotating superfluid Heliums\cite{Tkachenko1966}. In such a rotating dilute gas, the competition between inertial force, trapping potential and mean-field interaction results in a rich profile of vortex lattice, such as distinct core and global envelopes between mean-field quantum Hall and Thomas-Fermi regimes\cite{Fischer2003,Baym2004} and structural transitions to other lattice configurations in two-species bosons\cite{Mueller_Ho2002,Ueda2003,Cornell2}, dipolar systems\cite{Cooper2,HuiZhai1,Kumar2017}, Bose-Fermi double superfluids\cite{HuiZhai2}, etc.  %For instance, as increasing the boson density the vortex lattice can change from the lowest-Landau-level (LLL) Gaussian envelope to the Thomas-Fermi distribution showing an inverted parabola, with the core size continuously shrinking\cite{Fischer2003,Baym2004}. 
Despite these interesting phenomena, in practice it is quite  challenging  to explore the vortex lattice in rapidly rotating quantum Hall regime. A key obstacle is the exceedingly expanded cloud due to the cancellation between centrifugal and confinement forces ($\Omega=\omega$). This leads to the melting of vortex lattice with rather low density  and invisible pattern as observed in experiments\cite{QH_expt1, QH_expt2}. Alternatively, the boson density could be enhanced in the Landau-gauged quantum Hall regime by %geometrically 
squeezing the system with anisotropic traps\cite{MIT_expt1, MIT_expt2}, which  displays a crystalline vortex street\cite{Gora} instead of a triangular lattice as in symmetric case. 
 
  
Here, we show  above obstacle can be overcome in the quantum droplet of ultracold bosons. Such ultracold droplet, as balanced by the mean-field attraction and Lee-Huang-Yang(LHY) repulsion from quantum fluctuations\cite{Petrov2015}, has been successfully realized in both dipolar gases\cite{Ferrier-Barbut2016,Schmitt2016,Chomaz2016, Modugno,Pfau_4,Ferlaino_2} and alkali Bose-Bose mixtures\cite{Cabrera2018,Cheiney2018,Semeghini2018,DErrico2019,  DJWang}. Importantly, these droplets are self-bound and can be stabilized without any trap. They are thus expected to be immune from expansion as in previous cases of rotating Bose gases\cite{QH_expt1, QH_expt2}.
Although there have been several studies on a few vortices in dipolar droplets\cite{Cidrim2018:VorticesSelfboundDipolar,Roccuzzo2020:RotatingSupersolidDipolar,Ancilotto2021:VortexPropertiesExtended,Klaus2022:ObservationVorticesVortexa,Li2023:AnisotropicVortexQuantum,Bland2023:VorticesDipolarBoseEinstein} and binary boson droplets\cite{Li2018:TwodimensionalVortexQuantum,Zhang2019:SemidiscreteQuantumDroplets,Tengstrand2019:RotatingBinaryBoseEinstein,Caldara2022:VorticesQuantumDroplets,Kuan2023:DynamicsRapidlyRotating,Nikolaou2023:NovelSuperfluidStates}, the property of a rapidly rotating droplet in quantum Hall limit is still unknown at the moment. In particular, how will the self-bound nature and beyond-mean-field effect influence the vortex lattice is an important question to address.


% Fig.1: schematics of the setup and PV curve 
\begin{figure}[t]
    \centering
    \includegraphics[width=8.5cm]{fig1.pdf}
    \caption{Two limiting cases of vortex lattice for a 2D quantum droplet at different $\Omega/U$. Here $\Omega$ is the rotation frequency ($=$ trapping frequency $\omega$) and $U$ measures the interaction strength. For small $\Omega/U$ (a), each unit cell exhibits a rather small vortex core surrounded by a large flat-top area, resembling a needled surface. Increasing $\Omega/U$ to case (b), the system is frozen at the lowest-Landau-level, where the vortex core displays an extended Gaussian profile and the surface area becomes isolated and small. The dashed hexagon in each $n(x,y)$ plot marks a unit cell with characteristic length $l_{\Omega}=1/\sqrt{m\Omega}$. 
}\label{fig_schematic}
\end{figure}


In this work, we study the vortex lattice in a harmonically trapped two-dimensional (2D) Bose-Bose droplet with extreme rotation $\Omega=\omega$. Our calculations confirm that % the expectation that 
the self-bound nature of quantum droplet indeed protects the vortex lattice from melting, thus enabling a clear visualization of  triangular lattice in this limit. As shown in Fig.\ref{fig_schematic}, the exact vortex structure in a unit cell closely depends on the relative strength between $\Omega$ (or $\omega$) and interaction strength $U$. Namely, at $\Omega\ll U$ (Fig.\ref{fig_schematic}(a)), each vortex holds a rather small core with a large flat-top area around, such that the entire system resembles a needled surface. Increasing $\Omega/U$ will finally drive the system to the lowest-Landau-level (LLL), where the vortex core displays an extended Gaussian profile across the whole cell and the surface area becomes isolated and small (Fig.\ref{fig_schematic}(b)).
We emphasize that the flat-top surface here %, which does not appear in previous rotating Bose gases, 
inherits from the static droplet in vacuum and  reflects the crucial role played by LHY corrections. %This is in contrast to the usual vortex lattice in repulsive Bose gas, which displays  a varying global envelope instead of a flat one. 
Interestingly, the surface density here is always higher than that of a static droplet, and their ratio universally relies on a single parameter $\Omega/U$. These results have been demonstrated both numerically and from a variational approach. They will hopefully serve as a guideline for  future
experimental exploration of rapidly rotating ultracold droplets in quantum Hall limit.


We start from the energy functional of two-species bosons in a 2D rotating trap ($\hbar=1$):
\begin{eqnarray}
\cal{E}(\bm{\rho})&=&-\sum_{i=1,2}  \phi_i^{*}(\bm{\rho})\left( -\frac{\nabla^2}{2m} + \frac{m\omega^2 \bm{\rho}^2}{2} - \Omega L_z\right) \phi_i(\bm{\rho}) \nonumber\\
&&+ \sum_{ij}\frac{g_{ij}}{2}n_i(\bm{\rho})n_j(\bm{\rho})+ {\cal E}_{\rm LHY}(n_1(\bm{\rho}),n_2(\bm{\rho}));\label{E}
% &&+ \sum_{ij}\frac{g_{ij}}{2}n_i(\bm{\rho})n_j(\bm{\rho})+ {\cal E}_{\rm LHY}(n_i(\bm{\rho})); \label{E}
\end{eqnarray}
here $\bm{\bm{\rho}}=(x,y)$ is the 2D coordinate; $\phi_i$ is the wavefunction of the $i$-th species and $n_i=|\phi_i|^2$ is its density; $g_{ij}=4\pi a_{ij}/(ml_z)$ is the mean-field coupling between $i$- and $j$-species ($a_{ij}$ is the scattering length and $l_z$ is the confinement length along $z$ to generate 2D geometry); % ${\cal E}_{\rm LHY}$ is the LHY correction; 
$L_z=-i \bm{\rho}\times \nabla $ is the angular momentum along $z$. In this work we shall consider bosons near the mean-field collapse, i.e., with $\delta g\equiv g_{12}+ \sqrt{g_{11}g_{22}}\sim 0$, where the two-species densities are locked as $n_1/n_2=\sqrt{g_{22}/g_{11}}$\cite{Petrov2015}. In this case, one can employ the single-mode wave-function 
\begin{equation}
\Psi(\bm{\rho})=\left( \frac{\sqrt{g_{11}}+\sqrt{g_{22}}}{\sqrt{g_{22}}}\right)^{1/2} \phi_1 = \left( \frac{\sqrt{g_{11}}+\sqrt{g_{22}}}{\sqrt{g_{11}}}\right)^{1/2} \phi_2,
\end{equation}
and simplify (\ref{E}) as
\begin{eqnarray}
\cal{E}(\bm{\rho})&=&\Psi^{*}\left( -\frac{\nabla^2}{2m} + \frac{m\omega^2 \bm{\rho}^2}{2} - \Omega L_z\right)  \Psi + \frac{g}{2}n^2+ {\cal E}_{\rm LHY}(n). \nonumber\\ \label{E2}
\end{eqnarray}
Here  $n=|\Psi|^2$ is the total density, the reduced mean-field coupling is $g=2\delta g \sqrt{g_{11}g_{22}}/(\sqrt{g_{11}}+\sqrt{g_{22}})^2$, and the LHY energy reads ${\cal E}_{\rm LHY}=an^2 \ln(bn)$, with $a=2\pi a_{11}a_{22}/(ml_z^2)$ and $b=4\pi l_z \sqrt{ea_{11}a_{22}}$\citep{Petrov_2, Zin}. The ground state  can then be  obtained by imaginary-time evolution of the Gross-Pitaevskii(GP) equation $i\partial_t\Psi=\partial {\cal{E}}/\partial \Psi ^*$ based on (\ref{E2}). 

For a concrete example we shall take the two  hyperfine states of ${}^{39}$K atoms, $\ket{1} = \ket{F=1,m=-1}$ and $\ket{2}=\ket{F=1,m=0}$, as have been well studied in droplet experiments \cite{Semeghini2018,Cabrera2018,Cheiney2018}. Then we have $a_{11}=35a_B$ and $a_{12}=-53a_B$ ($a_B$ is the Bohr radius) and  $a_{22}$ is highly tunable through Feshbach resonance. Moreover, we assume a tunable $\Omega(=\omega)$ as well, while fix $l_z$ as $0.05\mu$m. % $2(2\pi)$Hz to $2(2\pi)$kHz, and $l_z$ is chosen as $0.05\mu$m. 
%The validity of single-mode treatment in (\ref{E2}) has been verified in our numerics.
In our numerics, a large boson number $N>10^4$ is considered, such that the droplet stays in thermodynamic limit with little finite-size effect. The boson filling factor $\nu=N/N_v$, with $N_v$ the vortex number, is $> 150$, which in turn validates the matter wave treatment in GP equation.  To facilitate later discussions, we introduce  a length scale $l_{\Omega}=1/\sqrt{m\Omega}$ to characterize the size of a unit cell, and an energy scale $U=an_0$ to characterize typical interaction strength, where  $n_0$ is the equilibrium density of a static droplet in vacuum. 
% near mean-field collapse. 

We have numerically obtained the ground state of rapidly rotating ${}^{39}$K droplets at $\Omega=\omega$.  For all values of $\Omega/U$, we find a stable and  regularly distributed triangular vortex lattice without melting, in contrast to those in repulsive Bose gases in the same quantum Hall limit\cite{QH_expt1, QH_expt2}. The vortex lattice is thus intrinsically self-bound, which will greatly facilitate its detection in experiments. Its internal structure, as summarized in Fig.\ref{fig_schematic},  falls into two limiting regimes depending on the ratio $\Omega/U$:

{\bf I. Needling regime:} when $\Omega\ll U$, the vortex core is very small as compared to the size of unit cell, and the rest large area is filled with bosons with flat-top (equilibrium) density. This structure, as shown in Fig.\ref{fig_schematic}(a), resembles a needled surface and thus it is called the needling regime.

{\bf II. LLL regime:} when $\Omega\gtrsim U$, the vortex spreads to the whole cell and the surface becomes isolated and individually small, see Fig.\ref{fig_schematic}(b). As shown later, the vortex profile in this case saturates at an extended Gaussian (Fig.\ref{fig_LLL}), a characteristic feature of the LLL regime.

%analytical analysis:
To physically understand these different structures, we write down the wave-function of the rotating droplet:
\begin{equation} 
    \Psi(\bm{\rho})=e^{i \Phi(\bm{\rho})}f(\bm{\rho})\sqrt{\bar{n}}, \label{wf}
\end{equation}
where $\Phi(\bm{\rho})$ is the phase, $f(\bm{\rho})$ is a real function rapidly varying on the scale of vortex core, and $\bar{n}$ is the equilibrium density far from the core. Since $f(\bm{\rho})$ has the same discrete translational symmetry as the vortex lattice, we can divide the lattice into Wigner-Seitz cells and treat $f$ in each cell separately. Then the vortex number is $N_v=\sum_j$ ($j$ is the cell index), and the total boson number is $N= N_v \nu$, with $\nu=\bar{n} \int_j d\bm{\rho} f^2$  the filling factor (here $\int_j$ is to integrate within a unit cell).

Now we aim to  express the total energy $E=\int d\bm{\rho} \cal{E}(\bm{\rho})$  in terms of $f$ and $\bar{n}$ in (\ref{wf}).  Since each vortex is singly quantized, we introduce the velocity field as $\bm{v}(\bm{\rho})\equiv \nabla \Phi/m$, which is contributed from an overall rigid rotation and a periodic local velocity, i.e., $\bm{v}(\bm{\rho})=\Omega \bm{e}_z\times \bm{\rho} + \bm{v}_l(\bm{\rho})$\cite{Tkachenko1966}. After straightforward algebra, we simplify the non-interacting part of $E$ as
%$E_0=N\Omega \alpha$, with
\begin{eqnarray}
E_0&=&N\Omega \alpha; \label{E0} \\
{\rm with} \ \    \alpha &=& \frac{ \int_j d\bm{\rho} \left( \frac{1}{2m}\big(\frac{\partial
     f}{\partial \rho}\big)^2 + \frac{m}{2} f^2 \bm{v}_l^2\right)}{ \Omega \int_j d\bm{\rho}f^2}. \nonumber
\end{eqnarray}
To exactly integrate  $\bm{v}_l^2$ in $\alpha$, we adopt a complex function $v_l(z) = i (\zeta^*(z)/m - \Omega z)$ in terms of $z=x+iy$, where $\zeta(z)$ is the Weierstrass Zeta function
$
    \zeta(z) = 1/{z} + \sum_{j\ne0}\big[{1/(z-z_j)} + 1/{z_j} + {z}/{z_j^2}\big],
$
and $\{z_j\}$ are complex coordinates of lattice sites \cite{Akhiezer1990:ElementsTheoryElliptic}. %The factor $\alpha$ depends on $\eta$ and the lattice symmetry, where the latter relates to the integral interval and the velocity field $v_{\rm r}$. 
Clearly $\alpha$ depends on the symmetry of vortex lattice, and its minimization  leads to the emergence of triangular lattice as shown later.

The interaction energy, as composed by mean-field and LHY parts, are given by
\begin{eqnarray}
&&\ \ \ \ \ \   E_{\rm int}=N\beta\Big(g\bar{n}/2 + a \bar{n}\ln(b\bar{n})\Big) + N\gamma a \bar{n}; \label{Eint} \\
&&{\rm with}  \ \ \ \ \beta = \dfrac{\int_j d\bm{\rho}f^4}{{\int_j d\bm{\rho}\, f^2}}, \ \ \ \  \gamma = \dfrac{\int_j d\bm{\rho}f^4\ln f^2}{{\int_j d\bm{\rho}\, f^2}}. \nonumber
\end{eqnarray}
The total energy per particle is then $\epsilon=(E_0+E_{\rm int})/N$:
\begin{equation}
\epsilon=\alpha\Omega+ \beta\big(g\bar{n}/2 + a \bar{n}\ln(b\bar{n})\Big) + \gamma a \bar{n}. \label{total_E}
\end{equation}
To this end, we have obtained $\epsilon$ as a function of $\bar{n}$ and three dimensionless quantities $\alpha,\ \beta,\ \gamma$ that are uniquely determined by $f$-function. Here we remark that the inclusion of LHY energy in $\epsilon$ is essential to support a self-bound solution at zero pressure, which is equivalent to requiring $\partial \epsilon/\partial \bar{n}=0$ and gives 
\begin{eqnarray}
    \bar{n}&=&e^{-\gamma/\beta}n_0, \label{nbar} \\
    \epsilon &=& \alpha\Omega + \beta e^{-\gamma/\beta} \epsilon_0.  \label{epsilon}
\end{eqnarray}
Here $n_0=b^{-1}e^{-1-g/(2a)}$ and $\epsilon_0=-an_0$ are, respectively, the equilibrium density and energy per particle for a static droplet in vacuum. 

To further parametrize $\alpha,\ \beta$ and $\gamma$, we now introduce a variational ansatz for $f(\rho)$ ($\rho=|\bm{\rho}|$) within a unit cell. Considering its definition in (\ref{wf}) with boundary condition $f(\rho\rightarrow 0)\propto \rho$, we write down the ansatz as 
\begin{equation}\label{trial_f}
    f({\rho})=\begin{cases}
        (\rho/l_c)\ e^{1-\rho/l_c}, & \rho\le l_c, \\
        1, & \rho> l_c.
    \end{cases} 
\end{equation}
Here $l_c$ characterizes the core radius, and one can check that both $f$ and $f'$ are continuous at $\rho=l_c$. Define a dimensionless quantity $\eta=l_c^2/l_{\Omega}^2$, which evaluates the core area with respect to the area of a unit cell, we can then parametrize $\alpha,\beta,\gamma$ solely by $\eta$. The ground state is then given by $\partial\epsilon /\partial \eta=0$, which leads to
\begin{equation}
  \dfrac{\Omega}{an_0}  \dfrac{\partial \alpha}{\partial \eta} =  \dfrac{\partial(\beta e^{-\gamma/\beta})}{\partial \eta}. \label{eta_eq}
\end{equation}
Importantly, Eq.(\ref{eta_eq}) implies that the solution of $\eta$, and thus other dimensionless quantities including  $\{\alpha,\beta,\gamma\}$ and the ratio $\bar{n}/n_0$, all rely on a single parameter $\Omega/(an_0)$. This parameter, according to Eq.(\ref{epsilon}), measures the relative strength between non-interaction and interaction energy scales. 


\begin{figure}[!ht]
    \centering
    \includegraphics[width=8cm]{fig2.pdf}
    \caption{Difference of energy per particle between the square and triangular vortex lattices $\Delta \epsilon= \epsilon_{\square}-\epsilon_{\bigtriangleup}$. The contributions from the non-interaction and interaction parts are  shown as $\Delta \epsilon_0$ and $\Delta \epsilon_{\rm int}$, respectively. The horizontal dashed line marks the value obtained in incompressible limit \cite{Tkachenko1966}. The energy unit here is the $\Omega$. }
    \label{fig_triangle}
\end{figure}


The above variational approach can well explain the structure of vortex lattices as shown in  Fig.\ref{fig_schematic}. First, it  gives the triangular lattice as the ground state configuration for all values of $\Omega/(an_0)$. To see this, in Fig.\ref{fig_triangle} we have shown the energy difference between the square and triangular lattices based on Eqs.(\ref{epsilon},\ref{eta_eq}). One can see that for all  $\Omega/(an_0)$, the triangular lattice is always more energetically favorable, and the energy gain is mostly contributed from the non-interacting part $\Delta \epsilon_0$, or equivalently, the difference in $\alpha$. In the incompressible limit with $\Omega/(an_0),\eta\rightarrow 0$, our calculation reproduces Tkachenko's result with $\frac{m^2}{2\pi}\left(\int_{\square}v_{l}^2-\int_{\bigtriangleup }v_{l}^2 \right)\approx 0.0105 $\cite{Tkachenko1966}.   

 

\begin{figure}[!ht]
    \centering
    \includegraphics[width=\linewidth]{fig3.pdf}
    \caption{Vortex lattice in the needling regime with $\Omega\ll an_0$. (a) Density profile $n(x,y=0)$  across three vortex cores at $\Omega/(an_0)=0.03$. Dashed line show the variational function fit according to Eq.(\ref{trial_f}) with optimized core radius $l_c$. %(b) $l_c/l_{\Omega}$ as a function of $\Omega/(an_0)$, where $l_{\Omega}=1/\sqrt{m\Omega}$ denotes the typical size of unit cell.  
    (b) $\bar{n}/n_0$ as a function of $\Omega/(an_0)$, with $\bar{n}$ (or $n_0$) the equilibrium surface density of rotating droplet (or static droplet in vacuum). Here $\Omega/(an_0)$ is changed  by varying different parameters ($\Omega$ or $\delta g$) and all data collapse into a single  curve (gray solid line),  signifying a {\it universal} dependence of $\bar{n}/n_0$ on $\Omega/(an_0)$. Black dashed line shows the results of variational approach based on Eq.(\ref{trial_f}), and the according $l_c/l_{\Omega}$ is shown in the inset plot.   }
    \label{fig_needling}
\end{figure}

Secondly, Eqs.(\ref{total_E}-\ref{eta_eq}) well explain the change of vortex core structure as tuning $\Omega/(an_0)$.
Let us start from the needling regime at $\Omega\ll an_0$. In this case, since the interaction part dominates in $\epsilon$, the system is largely unperturbed with an equilibrium density $\bar{n}$, while just a small region in each cell is significantly deformed by the vortex. In Fig.\ref{fig_needling}(a), we plot out a typical density profile for this case and show that it can be well fit using the variational ansatz in (\ref{trial_f}).
In Fig.\ref{fig_needling}(b), we further extract $\bar{n}/n_0$ as a function of $\Omega/(an_0)$. The data show that $\bar{n}$ for a rotating droplet is {\it always} higher than $n_0$ for a static one, and their ratio $\bar{n}/n_0$ gradually increases with $\Omega/(an_0)$. This can be attributed to the factor $e^{-\gamma/\beta}>1$ in (\ref{nbar}), given $\gamma<0$ according to its definition. The  variational approach is found to provide a quantitatively good prediction to $\bar{n}/n_0$ for small $\Omega/(an_0)\lesssim 0.15$, and it also predicts an increasing core size in each unit cell, see $l_c/l_{\Omega}\sim \Omega/(an_0)$ in the inset of Fig.\ref{fig_needling}(b). Remarkably, by choosing different variables ($\Omega$ or $\delta g$) in changing $\Omega/(an_0)$, we find that all data of $\bar{n}/n_0$ collapse into a single curve in Fig.\ref{fig_needling}(b). This demonstrates a {\it universal} dependence of $\bar{n}/n_0$ on $\Omega/(an_0)$, as suggested previously in the variational approach (see Eq.(\ref{eta_eq})). 

For $\Omega/(an_0)\gtrsim 0.15$,  the system no longer stays in the needling regime, as indicated by a considerably large core in a unit cell ($l_c/l_{\Omega}\sim 1$) and the breakdown of variational ansatz in predicting $\bar{n}/n_0$, see Fig.\ref{fig_needling}(b). In this case, the system gradually evolves to the LLL regime, and the energies in both Eqs.(\ref{E2},\ref{total_E}) are dominated by the non-interacting parts, i.e., with $\epsilon\sim \epsilon_0$ approaching the LLL energy $\Omega$ (thus $\alpha\rightarrow 1$). The associated $f(\rho)$ then develops an extended Gaussian distribution 
\begin{equation}
f(\rho)\rightarrow \rho e^{-\rho^2/(2l_{\Omega}^2)}. \label{Gaussian}
\end{equation}
Fig.\ref{fig_LLL}(a) shows that above function indeed determines the actual density within a unit cell at a large $\Omega/(an_0)=1$.


\begin{figure}[!ht]
    \centering
    \includegraphics[width=\linewidth]{fig4.pdf}
    \caption{Crossover of vortex lattice from the needling to the lowest-Landau-level(LLL) regime. (a) Same as Fig.\ref{fig_needling}(a) except for  $\Omega/(an_0)=1$. Dashed line show the fit to Gaussian function (\ref{Gaussian}). (b)  $\alpha=E_0/(N\Omega)$ as a function of $\Omega/(an_0)$.  (c) $y=\ln(|\Psi|/\rho)$ as a function of $x=\rho^2/l_{\Omega}^2$ for different $\Omega/(an_0)=0.02,\ 0.15$ and $0.92$, as marked by '(1),(2),(3)' in (b). For case (3), all points collapse into a straight line $y=-x/2+c$ (see dashed fit), which justifies the LLL wavefunction in (\ref{Gaussian}).  }
    \label{fig_LLL}
\end{figure}

The gradual crossover from the needling to LLL regimes is shown in Fig.\ref{fig_LLL}(b,c), where we have extracted $\alpha=E_0/(N\Omega)$ and $|\Psi|$ from exact numerics and plotted them as functions of $\Omega/(an_0)$. We can see that as increasing $\Omega/(an_0)$, $\alpha$ continuously decreases from a large value  to  $\sim 1$, suggesting the change of boson occupation from many higher Landau levels to the lowest one. During this process the vortex core function gradually develop the Gaussian form as in (\ref{Gaussian}), i.e., $\ln(|\Psi|/\rho)$ linearly depends on $\rho^2/l_{\Omega}^2$ with slope $-1/2$, see Fig.\ref{fig_LLL}(c). 

  

In summary, we have demonstrated a self-bound and clearly visible triangular vortex lattice in a rapidly rotating 2D Bose-Bose droplet in the quantum Hall limit. This is  in distinct contrast to the melted vortex lattice in the gaseous state of rotating bosons in the same limit. The LHY correction from quantum fluctuations plays an essential role in stabilizing such a rotating droplet. The revealed distinct structures of vortex lattice (see Fig.\ref{fig_schematic}) and the smooth crossover in-between, as well as the flat-top surface and its universal dependence on  $\Omega/U$ (Fig.\ref{fig_needling}(b))  can be readily detected in the current experiments of ultracold droplets. Our results suggest  the quantum droplet as a fascinating platform for the realistic exploration of quantum Hall regime with rotating bosons. In future, it will be interesting to investigate how the strongly correlated quantum Hall states emerge in this platform when the filling factor $\nu$ gets sufficiently low, as previously studied for single-species bosons\cite{Sinova2002:QuantumMeltingAbsence,Regnault2003:QuantumHallFractions}.
%, which  hopefully can stimulate more studies along this direction in future.


%Moreover, increasing the rotating frequency with respect to the interaction strength, the vortex lattice is found to undergo a smooth crossover from the needling regime to the LLL regime. Various properties, such as on the vortex core and the flat-top surface far from it, are studied both numerically and under a variational approach, which show quantitatively good consistence. Our results can be directly tested in the current experiments of ultracold droplets, and they also pave the way  for future exploration of bosonic quantum Hall regime in the intriguing self-bound objects. 



\acknowledgments
{\bf Acknowledgment.} We thank Hui Zhai for valuable discussions.
The work is supported by the National Key Research and Development Program of China (2018YFA0307600), the National Natural Science Foundation of China (12074419, 12134015), and the Strategic Priority Research Program of Chinese Academy of Sciences (XDB33000000).



\begin{thebibliography}{99}

\bibitem{review1} A. L. Fetter and A. A. Svidzinsky, \textit{Vortices in a trapped dilute Bose-Einstein condensate}, {J. Phys.: Condens. Matter {\bf 13}, R135 (2001)}.

 \bibitem{review2} N. R. Cooper, \textit{Rapidly rotating atomic gases}, {Advances in Physics {\bf 57}, 539 (2008)}.

 \bibitem{JasonHo} T.-L. Ho, \textit{Bose-Einstein Condensates with Large Number of Vortices}, {Phys. Rev. Lett. {\bf 87}, 060403 (2001)}.

\bibitem{Cooper} N. R. Cooper, N. K. Wilkin, and J. M. F. Gunn, \textit{Quantum Phases of Vortices in Rotating Bose-Einstein Condensates}, {Phys. Rev. Lett. {\bf 87}, 120405 (2001)}.


\bibitem{Abrikosov1957} A. A. Abrikosov, \textit{On the magnetic properties of superconductors of the second group}, {Sov. Phys. JETP {\bf 5}, 1174 (1957)}.
    
\bibitem{Tkachenko1966} V. K. Tkachenko, \textit{On vortex lattices}, {Sov. Phys. JETP {\bf 22}, 1282 (1966)}.

\bibitem{Fischer2003} U. R. Fischer and G. Baym, \textit{Vortex States of Rapidly Rotating Dilute Bose-Einstein Condensates},  {Phys. Rev. Lett. {\bf 90}, 140402 (2003)}.

\bibitem{Baym2004} G. Baym and C. J. Pethick, \textit{Vortex core structure and global properties of rapidly rotating Bose-Einstein condensates},  {Phys. Rev. A {\bf 69}, 043619 (2004)}.

\bibitem{Mueller_Ho2002} E. J. Mueller and T.-L. Ho, \textit{Two-Component Bose-Einstein Condensates with a Large Number of Vortices}, {Phys. Rev. Lett. {\bf 88}, 180403 (2002)}.

\bibitem{Ueda2003} K. Kasamatsu, M. Tsubota, and M. Ueda, \textit{Vortex Phase Diagram in Rotating Two-Component Bose-Einstein Condensates}, {Phys. Rev. Lett. {\bf 91}, 150406 (2003)}.

\bibitem{Cornell2} V. Schweikhard, I. Coddington, P. Engels, S. Tung, and A. E. A. Cornell, \textit{Vortex-Lattice Dynamics in Rotating Spinor Bose-Einstein Condensates},  {Phys. Rev. Lett. {\bf 93}, 210403 (2004)}.


\bibitem{Cooper2} N. R. Cooper, E. H. Rezayi, and S. H. Simon, \textit{Vortex Lattices in Rotating Atomic Bose Gases with Dipolar Interactions}, {Phys. Rev. Lett. {\bf 95}, 200402 (2005)}.

\bibitem{HuiZhai1} J. Zhang and H. Zhai, \textit{Vortex Lattices in Planar Bose-Einstein Condensates with Dipolar Interactions}, {Phys. Rev. Lett. {\bf 95}, 200403 (2005)}.

\bibitem{Kumar2017} R. K. Kumar, L. Tomio, B. A. Malomed, and A. Gammal, \textit{Vortex lattices in binary Bose-Einstein condensates with dipole-dipole interactions}, {Phys. Rev. A {\bf 96}, 063624 (2017)}.

\bibitem{HuiZhai2} Y. Jiang, R. Qi, Z.-Y. Shi, and H. Zhai, \textit{Vortex Lattices in the Bose-Fermi Superfluid Mixture}, {Phys. Rev. Lett. {\bf 118}, 080403 (2017)}.



 \bibitem{QH_expt1} V. Schweikhard, I. Coddington, P. Engels, V. P. Mogendorff, and E. A. Cornell, \textit{Rapidly Rotating Bose-Einstein Condensates in and near the Lowest Landau Level}, {Phys. Rev. Lett. {\bf 92}, 040404 (2004)}.

\bibitem{QH_expt2} V. Bretin, S. Stock, Y. Seurin, and J. Dalibard, \textit{Fast Rotation of a Bose-Einstein Condensate},  {Phys. Rev. Lett. {\bf 92}, 050403 (2004)}.

 
 \bibitem{MIT_expt1} R. J. Fletcher, A. Shaffer, C. C. Wilson, P. B. Patel, Z. Yan, V. Crépel, B. Mukherjee, and M. W. Zwierlein, \textit{Geometric squeezing into the lowest Landau level},  {Science {\bf 372}, 1318 (2021)}.

 
 \bibitem{MIT_expt2} B. Mukherjee, A. Shaffer, P. B. Patel, Z. Yan, C. C. Wilson, V. Crépel, R. J. Fletcher, and M. Zwierlein, \textit{Crystallization of bosonic quantum Hall states in a rotating quantum gas},  {Nature {\bf 601}, 7891 (2022)}.

 
 \bibitem{Gora} S. Sinha and G. V. Shlyapnikov, \textit{Two-Dimensional Bose-Einstein Condensate under Extreme Rotation}, {Phys. Rev. Lett. {\bf 94}, 150401 (2005)}.


\bibitem{Petrov2015} D. S. Petrov, \textit{Quantum Mechanical Stabilization of a Collapsing Bose-Bose Mixture},  {Phys. Rev. Lett. {\bf 115}, 155302 (2015)}.


%droplet expt:

\bibitem{Ferrier-Barbut2016} I. Ferrier-Barbut, H. Kadau, M. Schmitt, M. Wenzel, and T. Pfau, \textit{Observation of Quantum Droplets in a Strongly Dipolar Bose Gas}, Phys. Rev. Lett. {\bf 116}, 215301 (2016).

\bibitem{Schmitt2016} M. Schmitt, M. Wenzel, F. B{\"o}ttcher, I. Ferrier-Barbut, and T. Pfau, \textit{Self-bound droplets of a dilute magnetic quantum liquid}, Nature {\bf 539}, 259 (2016).

\bibitem{Chomaz2016} L. Chomaz, S. Baier, D. Petter, M. J. Mark, F. W{\"a}chtler, L. Santos, and F. Ferlaino, \textit{Quantum-Fluctuation-Driven Crossover from a Dilute Bose-Einstein Condensate to a Macrodroplet in a Dipolar Quantum Fluid}, Phys. Rev. X {\bf 6}, 041039 (2016).

%supersolid in dipolar droplet:
\bibitem{Modugno}L. Tanzi, E. Lucioni, F. Fama, J. Catani, A. Fioretti, C. Gabbanini, R. N. Bisset, L. Santos, and G. Modugno, Observation of a Dipolar Quantum Gas with Metastable Supersolid Properties, Phys. Rev. Lett. {\bf 122}, 130405 (2019).
\bibitem{Pfau_4}F. B{\" o}ttcher, J.-N. Schmidt, M. Wenzel, J. Hertkorn, M. Guo, T. Langen, and T. Pfau, Transient Supersolid Properties in an Array of Dipolar Quantum Droplets, Phys. Rev. X {\bf 9}, 011051 (2019).
\bibitem{Ferlaino_2}L. Chomaz, D. Petter, P. Ilzh{\"{o}}fer, G. Natale, A. Trautmann, C. Politi, G. Durastante, R.M.W. van Bijnen, A. Patscheider, M. Sohmen, M.J. Mark, and F. Ferlaino, Long-Lived and Transient Supersolid Behaviors in Dipolar Quantum Gases, Phys. Rev. X {\bf 9}, 021012 (2019).


\bibitem{Cabrera2018} C. R. Cabrera, L. Tanzi, J. Sanz, B. Naylor, P. Thomas, P. Cheiney, and L. Tarruell, \textit{Quantum liquid droplets in a mixture of Bose-Einstein condensates}, Science {\bf 359}, 301 (2018).


\bibitem{Cheiney2018} P. Cheiney, C. R. Cabrera, J. Sanz, B. Naylor, L. Tanzi, and L. Tarruell, \textit{Bright Soliton to Quantum Droplet Transition in a Mixture of Bose-Einstein Condensates}, Phys. Rev. Lett. {\bf 120}, 135301 (2018).

\bibitem{Semeghini2018} G. Semeghini, G. Ferioli, L. Masi, C. Mazzinghi, L. Wolswijk, F. Minardi, M. Modugno, G. Modugno, M. Inguscio, and M. Fattori, \textit{Self-Bound Quantum Droplets of Atomic Mixtures in Free Space}, Phys. Rev. Lett. {\bf 120}, 235301 (2018).


\bibitem{DErrico2019} C. D'Errico, A. Burchianti, M. Prevedelli, L. Salasnich, F. Ancilotto, M. Modugno, F. Minardi, and C. Fort, \textit{Observation of quantum droplets in a heteronuclear bosonic mixture}, Phys. Rev. Res. {\bf 1}, 033155 (2019).

\bibitem{DJWang}Z. Guo, F. Jia, L. Li, Y. Ma, J. M. Hutson, X. Cui, and D. Wang, Lee-Huang-Yang effects in the ultracold mixture of $^{23}$Na and $^{87}$Rb with attractive interspecies interactions, Phys. Rev. Res. {\bf 3}, 033247 (2021).


    % few vortices in dipolar droplet
    \bibitem{Cidrim2018:VorticesSelfboundDipolar} A. Cidrim, F. E. A. dos Santos, E. A. L. Henn, and T. Macrì, \textit{Vortices in self-bound dipolar droplets},  {Phys. Rev. A {\bf 98}, 023618 (2018)}.
    \bibitem{Roccuzzo2020:RotatingSupersolidDipolar} S. M. Roccuzzo, A. Gallemí, A. Recati, and S. Stringari, \textit{Rotating a Supersolid Dipolar Gas},  {Phys. Rev. Lett. {\bf 124}, 045702 (2020)}.
    \bibitem{Ancilotto2021:VortexPropertiesExtended} F. Ancilotto, M. Barranco, M. Pi, and L. Reatto, \textit{Vortex properties in the extended supersolid phase of dipolar Bose-Einstein condensates},  {Phys. Rev. A {\bf 103}, 033314 (2021)}.
    \bibitem{Klaus2022:ObservationVorticesVortexa} L. Klaus, T. Bland, E. Poli, C. Politi, G. Lamporesi, E. Casotti, R. N. Bisset, M. J. Mark, and F. Ferlaino, \textit{Observation of vortices and vortex stripes in a dipolar condensate},  {Nat. Phys. {\bf 18}, 12 (2022)}.
    \bibitem{Li2023:AnisotropicVortexQuantum} G. Li, X. Jiang, B. Liu, Z. Chen, B. A. Malomed, and Y. Li, \textit{Anisotropic vortex quantum droplets in dipolar Bose-Einstein condensates}, arXiv:2301.04305.
    \bibitem{Bland2023:VorticesDipolarBoseEinstein} T. Bland, G. Lamporesi, M. J. Mark, and F. Ferlaino, \textit{Vortices in dipolar Bose-Einstein condensates}, arXiv:2303.13263.
    
    % few vortices in binary droplet
    \bibitem{Li2018:TwodimensionalVortexQuantum} Y. Li, Z. Chen, Z. Luo, C. Huang, H. Tan, W. Pang, and B. A. Malomed, \textit{Two-dimensional vortex quantum droplets},  {Phys. Rev. A {\bf 98}, 063602 (2018)}.
    \bibitem{Zhang2019:SemidiscreteQuantumDroplets} X. Zhang, X. Xu, Y. Zheng, Z. Chen, B. Liu, C. Huang, B. A. Malomed, and Y. Li, \textit{Semidiscrete quantum droplets and vortices},  {Phys. Rev. Lett. {\bf 123}, 133901 (2019)}.
    \bibitem{Tengstrand2019:RotatingBinaryBoseEinstein} M. N. Tengstrand, P. Stürmer, E. Karabulut, and S. M. Reimann, \textit{Rotating binary Bose-Einstein condensates and vortex clusters in quantum droplets},  {Phys. Rev. Lett. {\bf 123}, 160405 (2019)}.
    \bibitem{Caldara2022:VorticesQuantumDroplets} M. Caldara and F. Ancilotto, \textit{Vortices in quantum droplets of heteronuclear Bose mixtures},  {Phys. Rev. A {\bf 105}, 063328 (2022)}.
    \bibitem{Kuan2023:DynamicsRapidlyRotating} W.-H. Kuan, Y.-W. Wang, and S.-C. Cheng, \textit{Dynamics of Rapidly Rotating Bose-Einstein Quantum Droplets}, arXiv:2302.07481.
    \bibitem{Nikolaou2023:NovelSuperfluidStates} S. Nikolaou, G. M. Kavoulakis, and M. Ogren, \textit{Novel superfluid states in rotating quantum droplets confined in a harmonic potential}, arXiv:2305.09422.

%2D LHY:
\bibitem{Petrov_2}D. S. Petrov and G. E. Astrakharchik, Ultradilute Low-Dimensional Liquids,  Phys. Rev. Lett. {\bf 117}, 100401 (2016).
\bibitem{Zin}P. Zin, M. Pylak, T. Wasak, M. Gajda, and Z. Idziaszek, Quantum Bose-Bose droplets at a dimensional crossover, Phys. Rev. A {\bf 98}, 051603(R) (2018).
%\bibitem{Cui}

    % Weierstrass functions
    \bibitem{Akhiezer1990:ElementsTheoryElliptic} N. I. Akhiezer, \textit{Elements of the Theory of Elliptic Functions}, New Ed edition (American Mathematical Society, Providence, RI, 1990).



% quasi-2d GPE
%    \bibitem{Ma2023:QuantumfluctuationdrivenDynamicsDroplet} Y. Ma and X. Cui, \textit{Quantum-fluctuation-driven dynamics of droplet splashing, recoiling, and deposition in ultracold binary Bose gases},  {Phys. Rev. Res. {\bf 5}, 013100 (2023)}.

    \bibitem{Sinova2002:QuantumMeltingAbsence} J. Sinova, C. B. Hanna, and A. H. MacDonald, \textit{Quantum Melting and Absence of Bose-Einstein Condensation in Two-Dimensional Vortex Matter}, %\href{https://doi.org/10.1103/PhysRevLett.89.030403} 
    {Phys. Rev. Lett. {\bf 89}, 030403 (2002)}.
    \bibitem{Regnault2003:QuantumHallFractions} N. Regnault and Th. Jolicoeur, \textit{Quantum Hall Fractions in Rotating Bose-Einstein Condensates}, %\href{https://doi.org/10.1103/PhysRevLett.91.030402} 
    {Phys. Rev. Lett. {\bf 91}, 030402 (2003)}.
    

%\bibitem{supple} See supplementary materials for more details.

\end{thebibliography}

\end{document}
